\documentclass[a4paper, 12pt]{article}

% ===== 中文支持 =====
\usepackage[UTF8]{ctex}

% ===== 页面设置 =====
\usepackage[top=2.5cm, bottom=2.5cm, left=3cm, right=3cm]{geometry}
\usepackage{setspace}
\onehalfspacing

% ===== 数学环境 =====
\usepackage{amsmath, amssymb, amsfonts}
\usepackage{bm}

% ===== 表格 =====
\usepackage{booktabs}
\usepackage{longtable}
\usepackage{multirow}
\usepackage{array}

% ===== 图片 =====
\usepackage{graphicx}
\usepackage{subcaption}  % 提供 subfigure 环境（并排图），替代已废弃且与 caption 冲突的 subfigure 宏包

% ===== 算法 =====
\usepackage{algorithm}
\usepackage{algorithmicx}
\usepackage{algpseudocode}

% ===== 代码框 =====
% [most] 选项已自带 listings 与 xcolor 支持，无需额外引入
\usepackage[most]{tcolorbox}

% 可跨页的代码框
\newtcblisting{codebox}[1][]{%
  listing only,
  colback=gray!5,
  colframe=black!70,
  fontupper=\scriptsize\ttfamily,
  boxrule=0.8pt,
  arc=2pt,
  left=4pt, right=4pt, top=2pt, bottom=2pt,
  breakable,
  #1
}

% ===== 超链接 =====
\usepackage{hyperref}
\hypersetup{
  colorlinks=true,
  linkcolor=black,
  citecolor=blue,
  urlcolor=blue
}

% ===== 标题设置 =====
\usepackage{titlesec}
\titleformat{\section}{\Large\bfseries}{\thesection}{1em}{}
\titleformat{\subsection}{\large\bfseries}{\thesubsection}{1em}{}
\titleformat{\subsubsection}{\normalsize\bfseries}{\thesubsubsection}{1em}{}

% ===== 参考文献 =====
\usepackage[numbers]{natbib}

% ===== 其他 =====
\usepackage{enumitem}
\usepackage{float}
\usepackage{caption}
\captionsetup{font=small, labelfont=bf}

\title{\textbf{论文标题}}
\author{作者姓名 \\ 学校名称}
\date{\today}

\begin{document}

\maketitle

% ===== 摘要 =====
\begin{abstract}
\noindent
这里是摘要内容。开头段用 2-3 句话概括选题背景与核心研究目标。

针对问题一，本文采用...方法，解决了...问题，得到...结果。

针对问题二，本文采用...方法，解决了...问题，得到...结果。

针对问题三，本文采用...方法，解决了...问题，得到...结果。

本文的核心贡献在于...。

\textbf{关键词：}关键词1；关键词2；关键词3；关键词4
\label{abstract:end}
\end{abstract}

% 摘要单独占一页（编译后验证 abstract:end 页码=1）；正文章节之间不再人为断页
\newpage

% ===== 一、问题重述 =====
\section{问题重述}

\subsection{问题背景}
% 选题的现实意义，学术化语言改写

\subsection{核心子问题}
% 按赛题顺序提炼具体任务

% ===== 二、问题分析 =====
\section{问题分析}

\subsection{问题一分析}
% 相关条件 → 整体分析 → 初步方法
% 思维流程图引用
\begin{figure}[H]
  \centering
  \includegraphics[width=0.9\textwidth]{figures/problem1/flow.pdf}
  \caption{问题一分析思路流程图}
  \label{fig:flow1}
\end{figure}

\subsection{问题二分析}
% 同上

\subsection{问题三分析}
% 同上

% ===== 三、模型假设 =====
\section{模型假设}

为简化问题，本文做出以下假设：

\begin{itemize}
  \item \textbf{假设1：}题目明确给出的假设条件……
  \item \textbf{假设2：}排除小概率事件……
  \item \textbf{假设3：}仅考虑核心因素，忽略次要因素……
  \item \textbf{假设4：}模型本身要求的假设……
  \item \textbf{假设5：}对参数形式/分布的假设……
  \item \textbf{假设6：}用于简化的假设……
\end{itemize}

% ===== 四、符号说明 =====
\section{符号说明}

% 全论文表格统一 longtable + p{}，列宽比例总和 = 1.04 - 0.04×N（此处 3 列 = 0.92）
% 列序与 writing-standards.md §7 一致：符号 | 说明 | 单位
\begin{longtable}{>{\centering\arraybackslash}p{0.16\textwidth}>{\centering\arraybackslash}p{0.60\textwidth}>{\centering\arraybackslash}p{0.16\textwidth}}
  \caption{符号说明}
  \label{tab:symbols} \\
  \toprule
  符号 & 说明 & 单位 \\
  \midrule
  \endfirsthead
  \caption{符号说明（续）} \\
  \toprule
  符号 & 说明 & 单位 \\
  \midrule
  \endhead
  \bottomrule
  \endfoot
  $X_1$ & 第 1 区域日客流量 & 人 \\
  $T$ & 时间周期 & 小时 \\
  $\alpha$ & 学习率 & --- \\
\end{longtable}

% ===== 五、模型建立与求解 =====
\section{模型的建立与求解}

\subsection{问题一模型建立}
% 将实际问题抽象为数学表达式
% 核心公式居中展示
\begin{equation}
  \min_{x} f(x) = \sum_{i=1}^{n} c_i x_i
  \label{eq:obj1}
\end{equation}

约束条件：
\begin{equation}
  \begin{cases}
    \sum_{j=1}^{m} a_{ij} x_j \geq b_i, & i = 1, 2, \ldots, p \\
    x_j \geq 0, & j = 1, 2, \ldots, n
  \end{cases}
  \label{eq:cons1}
\end{equation}

% 公式解释文字

\subsection{问题一模型求解}
% 软件选择及理由
% 计算步骤
% 结果呈现

\begin{longtable}{>{\centering\arraybackslash}p{0.30\textwidth}>{\centering\arraybackslash}p{0.31\textwidth}>{\centering\arraybackslash}p{0.31\textwidth}}
  \caption{问题一求解结果}
  \label{tab:result1} \\
  \toprule
  方案 & 目标函数值 & 优化率 \\
  \midrule
  \endfirsthead
  \caption{问题一求解结果（续）} \\
  \toprule
  方案 & 目标函数值 & 优化率 \\
  \midrule
  \endhead
  \bottomrule
  \endfoot
  方案A & 1234.56 & --- \\
  方案B & 1089.23 & 11.8\% \\
\end{longtable}

\subsection{问题二模型建立}
% ...

\subsection{问题二模型求解}
% ...

% ===== 六、模型分析与检验 =====
\section{模型的分析与检验}

\subsection{灵敏度分析}

\begin{figure}[H]
  \centering
  \includegraphics[width=0.8\textwidth]{figures/sensitivity.png}
  \caption{灵敏度分析结果}
  \label{fig:sensitivity}
\end{figure}

\subsection{误差分析}

\subsection{模型检验}

% ===== 七、模型评价、改进与推广 =====
\section{模型的评价、改进与推广}

\subsection{模型优点}
\begin{itemize}
  \item \textbf{优点1：}……
  \item \textbf{优点2：}……
  \item \textbf{优点3：}……
  \item \textbf{优点4：}……
\end{itemize}

\subsection{模型缺点}
\begin{itemize}
  \item \textbf{缺点1：}……
  \item \textbf{缺点2：}……
\end{itemize}

\subsection{模型改进}
% 针对缺点的可落地方案

\subsection{模型推广}
% 扩展至类似问题

% ===== 参考文献 =====
\begin{thebibliography}{99}
  \bibitem{ref1} 作者. 文献标题[J]. 期刊名, 年份, 卷(期): 页码.
  \bibitem{ref2} 作者. 书名[M]. 出版地: 出版社, 年份.
\end{thebibliography}

% ===== 附录 =====
% 用 \appendix 命令（不是环境）切换到附录模式，其后 \section 自动编号为 A、B、C…
\appendix

\section{附件清单}

\begin{longtable}{>{\centering\arraybackslash}p{0.10\textwidth}>{\centering\arraybackslash}p{0.25\textwidth}>{\centering\arraybackslash}p{0.60\textwidth}}
  \caption{附件清单}
  \label{tab:appendix} \\
  \toprule
  编号 & 文件名 & 说明 \\
  \midrule
  \endfirsthead
  \caption{附件清单（续）} \\
  \toprule
  编号 & 文件名 & 说明 \\
  \midrule
  \endhead
  \bottomrule
  \endfoot
  % 根据赛题实际附件信息填写，代码文件格式可能为 .m / .py 等
  附件1 & xxx.xlsx & 题目给定的数据文件 \\
  附件2 & xxx.py / xxx.m & 代码主程序 \\
  附件3 & xxx.xlsx & 计算结果或中间数据 \\
\end{longtable}

\end{document}
